% 01_introduction.tex

\section{Introduction}
\label{sec:introduction}

In standard costly signaling, separation is sustained from below: low types cannot afford high-cost signaling, so they do not mimic high types. The model here studies the opposite failure mode: costly signaling can reverse when the relevant cost is falsification exposure rather than production effort. A good type may still benefit from an ambitious signal today, but can rationally retreat because that action creates a one-period exposure that can be repriced against her future credibility.

The mechanism is dynamic and belief-driven. The ambitious action has two effects. It is informative in the current period, but it also leaves a claim exposed to future scrutiny. If high scrutiny arrives tomorrow, the claim can be repriced downward and impose an additional penalty. This is adverse repricing, not a true-type audit.

The primitives are minimal: two types, two sender signals, and an exogenous binary judgment state. The sender observes the low-scrutiny risk $q=\Pr(j_{t+1}=H\mid j_t=L)$ before acting. If she sends the ambitious signal, tomorrow's state can impose an extra exposure cost. If she sends the pooling-safe action, that one-period shadow disappears.

The model has four moving parts.
\begin{enumerate}
\item A positioning section that explains why the reversal is not a trivial omission from the post-Spence literature.
\item A D1-closed fixed-point core that shows a top-down signaling reversal: high types exit separation first when $q$ is large, and D1 then removes the high-type meaning of off-path ambition.
\item An observable-implications section that translates the binary collapse into support polarization in a richer type space.
\item A dynamic forward-invariance result that adds trust-capital/semantic dynamics and, under a persistence law for $q_t$, shows endogenous accumulation does not automatically restore separation.
\end{enumerate}

The first stage explains why the reverse was not immediate: the friction changes from production cost to reputation exposure, the dynamic object is a state transition rather than gradual accumulation alone, and honesty becomes an endogenous asset rather than only an exogenous preference. The second stage identifies how adverse repricing plus a D1-style off-path map can select pooling from above. It closes the belief loop in binary form: a self-referential interpretation where high occupancy of the ambitious signal implies high-type meaning, but then current IC may already fail, and off-path reinterpretation then collapses that meaning. The third stage records the observable cross-sectional shadow. The final stage adds dynamic state variables and gives the defensive result: under a law that keeps $q_t$ above the collapse boundary, endogenous accumulation does not restore separation automatically.

The result is defensive: no new stopping phenomenon appears, and the paper does not claim that all dynamic systems remain collapsed. It shows that, under the persistence law, learning and accumulation alone do not force separation to return.

Section~\ref{sec:why-not-before} explains why the reverse was not a trivial omission. Section~\ref{sec:primal-model} introduces primitives and the one-period IC comparison. Section~\ref{sec:psychological-exposure} gives one primitive foundation for the reversed exposure-cost ranking \(\Phi_H>\Phi_L\). Section~\ref{sec:belief-loop} proves the D1-closed fixed-point theorem and the threshold $q^*=\max\{q_{\mathrm{IC}},q_{\mathrm{D1}}\}$. Section~\ref{sec:polarization} presents the observable polarization implication. Section~\ref{sec:dynamic-invariance} gives the dynamic forward-invariance result under the persistence law for $q_t$.
